
\begin{ledgroupsized}[r]{120mm}
\footnotesize
\pstart
\noindent\textbf{Überlieferung:}
\pend
\end{ledgroupsized}
\begin{ledgroupsized}[r]{114mm}
\footnotesize
\pstart
\parindent -6mm
\makebox[6mm][l]{\textit{L}}%
Konzept:
LH~XXXV~8, 30~Bl.~40.
Ein Blatt~2\textsuperscript{o}
(ca.~16~x~20~cm.)
{Wasserzeichen} in der Blattmitte;
oberer Rand nach Abschnitt unregelmäßig.
Zwei Seiten;
\pend
\end{ledgroupsized}
\vspace{5mm}
\begin{ledgroup}
\footnotesize
\pstart
\noindent
\textbf{Datierungsgründe:}
Das Wasserzeichen im Träger von N.\;\ref{42014}
konnte noch nicht identifiziert werden.
Die Datierung beruht auf bedeutenden inhaltlichen Ähnlichkeiten mit anderen Konzepten und mit dem Briefwechsel mit \indexN{\textls{Papin} (Papinus), Denis 1647\textendash1713}Denis Papin.
\pend
\pstart
Die Analyse komplexer Stoßfälle anhand zeitlicher Phasen ("status primus," "secundus" usw.) findet sich nicht nur im Ansatz, sondern auch in den Einzelheiten von Text und Zeichnungen in Konzepten wie N.\;\ref{41780} (Mitte September 1690 bis 1696), N.\;\ref{58300} (Januar 1691 bis 7. Februar 1697) und  N.\;\ref{58297} (August 1697 bis April 1703) wieder.
Dabei weist N.\;\ref{42014} eine engere Nähe zu beiden früheren Konzepten auf: Die Fragen werden hier aufgeworfen aber teilweise offen gelassen, oder auf eine später zu leistende Untersuchung vertagt (die auf 
\refnote{FnC:LH_35_08_30_040-advena}
nur skizzierte Fallstudie leistet Leibniz tatsächlich ausführlich im Konzept N.\;\ref{58300}).
Auf \refnote{FnC:LH_35_08_30_040-Papin} bezieht sich Leibniz auf das Gedankenexperiment einer zwischen den stoßenden Körpern befindlichen Feder mit den Worten "quae omnia distinctius exponi merentur".
Genau diesen Fall schlägt aber Leibniz im Juli 1696 Papin zur Diskussion vor; er wird beide Kontrahenten bis Januar 1698 rege beschäftigen (siehe dazu die Datierung zu N.\;\ref{60360} und die dort zitierten Briefe).
Die Ausdrucksweise von N.\;\ref{42014} lässt den Schluss zu, dass zu diesem Zeitpunkt die Auseinandersetzung mit Papin noch nicht stattgefunden (oder unlängst begonnen) hatte.
Aus den genannten Gründen ergibt sich das Jahr 1696 als Terminus ante quem.
\pend
\pstart
Als Terminus post quem wird das oben genannte Konzept N.\;\ref{41780} angenommen. Dort, wie auch im Schlussteil von N.\;\ref{42014} mit dem Diagramm \bracks{\textit{Fig. 4}}, hinterfragt (und verneint) Leibniz in ähnlichen Worten eine These zum schiefen Mehrkörperstoß, die Leibniz aus \indexN{\textls{Huygens} (Hugenius, Ugenius, Hugens, Huguens), Christiaan 1629\textendash1695}Huygens' 
\cite{00525}\textit{Traité de la lumière}, Leiden 1690, kennt (und in N.\;\ref{41780} ausdrücklich Huygens zuschreibt: \refnote{FnC:37_03_067_Huygens}).
Die Lektüre des Buchs kann frühestens Mitte September 1690 stattgefunden haben: Siehe die Datierung von N.\;\ref{41780} und des kurzen Exzerptes aus der einschlägigen Stelle des \textit{Traité}, das bei dieser Lektüre entstanden sein muss (in \cite{02094}\textit{LSB} VIII,~3 N.\;75).
\pend
\end{ledgroup}
\selectlanguage{latin}
\vspace{8mm}
\pstart
\centering
\edtext{\edtext{%
\href{\Scan{LH_35_08_30_040r}}%
{\lbrack40~r\textsuperscript{o}\rbrack} 
\hfill
\textls{Suppositio}}{%
\lemma{\href{\Scan{LH_35_08_30_040r}}%
{\lbrack40~r\textsuperscript{o}\rbrack}}%
\Bfootnote{%
\textit{(1)}~Sit recta
\textit{(2)}~\bars{In \textit{streicht Hrsg.}}
\textit{(a)}~recta
\textit{(b)}~linea signata sint puncta aequidistantia
\textit{(3)}~Suppositio: si globus durus directe occurrat 
\textit{(4)}~\textls{Suppositio}~\textit{L}}}
\textls{experimentis confirmata}\indexS{experimentum}}{%
\lemma{\textls{Suppositio experimentis confirmata}}%
\Cfootnote{Siehe bspw. \indexN{\textls{Mariotte}, Edme, Seigneur de Chazeuil ca. 1620\textendash1684}\textsc{E. Mariotte}, \cite{00311}\textit{Traité de la percussion}, Paris 1673, Première partie, Prop. XVI, S.\;100--102.}}
\hfill
\phantom{\lbrack40~r\textsuperscript{o}\rbrack}\ 
\pend
\nobreak
\vspace{.5em}
\pstart
\noindent
Si globus\indexS{globus}\indexS{globus incurrens} durus 
\edtext{incurrat directe}{%
\lemma{incurrat}%
\Bfootnote{%
\textit{(1)}~in ali
\textit{(2)}~directe~\textit{L}}}
et centraliter\indexS{incursus directus et centralis globi in aequalem quiescentem} in alium aequalem et materia similem sed 
\edtext{quiescentem;\indexS{globus quiescens} 
eveniet ut post ictum\indexS{ictus} incurrens\indexS{globus incurrens} quietem\indexS{quies} accipiat in loco}{%
\lemma{quiescentem;}%
\Bfootnote{%
\textit{(1)}~post ictum 
\textit{(a)}~quiescen
\textit{(b)}~incurrens quietem
\textit{(aa)}~ex
\textit{(bb)}~accipiet
\textit{(cc)}~accipiet
\textit{(2)}~con
\textit{(3)}~eveniet \lbrack...\rbrack\ accipiat
\textit{(a)}~loco e
\textit{(b)}~in loco~\textit{L}}}
excipientis;\indexS{globus excipiens} excipiens autem motum\indexS{motus} accipiat incurrentis, 
eademque celeritate et directione\indexS{celeritas et directio} qua ille sine incursu\indexS{incursus} perrexisset moveatur.
\pend
\begin{edfigure}[1em][1em]
\includegraphics[scale=1]{%
gesamttex/edit_VIII-4/images/LH_35_08_30_040_d1_040r.pdf}
\edcaption{\edtext{\lbrack\textit{Fig.~1}\rbrack}{%
\lemma{\lbrack\textit{Fig.~1}\rbrack}%
\Cfootnote[nonum]{Leibniz hat im oberen Bereich von Bl.~40~r\textsuperscript{o} das Diagramm in zwei gleichen Fassungen gezeichnet; die versehentlich stehen gebliebene Dublette, sowie ein gestrichener Entwurf zum Diagramm, werden nicht wiedergegeben.}}}
\end{edfigure}
\pstart
\edtext{%
Ut si in
\edtext{recta (fig.~1) signentur tria puncta aequidistantia et sit}{%
\lemma{recta}%
\Bfootnote{%
\textit{(1)}~in qua signentur tria puncta aequidistantia moto \textit{B}
\textit{(2)}~(fig.~1) signentur tria puncta aequidistantia
\textit{(a)}~sitque \textit{B}
\textit{(b)}~et sit~\textit{L}}}
\textit{B} in motu, \textit{L} vero quiescat ut in statu primo (fig.~1) inde incurrat \textit{B} in \textit{L} statu secundo, denique statu
\edtext{tertio quiesc\bracks{a}t \textit{B}}{%
\lemma{tertio}%
\Bfootnote{%
\textit{(1)}~quiescat
\textit{(2)}~\bars{quiescet \textit{ändert Hrsg.}} \textit{B}~\textit{L}}}
in loco ipsius 
\edtext{\textit{L}, nempe}{%
\lemma{\textit{L},}%
\Bfootnote{%
\textit{(1)}~sed
\textit{(2)}~nempe~\textit{L}}}
2, sed
\edtext{\lbrack\textit{L}\rbrack}{\lemma{}\Bfootnote{\textit{B} \textit{L ändert Hrsg.}}}
accipiat motum 2\,3 aequalem ipsi motui 1\,2 quem initio habebat \textit{B}.%
}{\lemma{}\Bfootnote{Ut si \lbrack...\rbrack\ habebat \textit{B}. \textit{erg.~L}}}
\pend
\vspace{0.5\baselineskip}
\pstart
\centering
\edlabel{35_08_30_040_1a}%
\edtext{}{%
\xxref{35_08_30_040_1a}{35_08_30_040_1b}%
\lemma{\textls{Theorema:}}%
\Bfootnote{%
\textit{(1)}~Si
\textit{(a)}~duo
\textit{(b)}~globi in
\textit{(2)}~\bars{Si \textit{streicht Hrsg.}} duo globi
\textit{(a)}~in re
\textit{(b)}~aequales m
\textit{(3)}~Si sint quatuor globi 
\textit{(a)}~in eadem recta
\textit{(b)}~duri~\textit{L}}}%
\textls{Theorema}:
\pend
\nobreak
\pstart
\noindent
Si sint quatuor globi\indexS{globi quatuor aequales} duri\edlabel{35_08_30_040_1b}
aequales et materia similes, centris\indexS{centrum gravitatis} in eadem recta positis, ex quibus duo moveantur, et duo
\edtext{quiescant, absoluto}{%
\lemma{quiescant,}%
\Bfootnote{%
\textit{(1)}~post
\textit{(2)}~absolutis
\textit{(3)}~absoluto~\textit{L}}}
concursu\indexS{concursus absolutus} fiet, ut duo incurrentes\indexS{globi incurrentes} loco excipientium\indexS{globi excipientes} quietem\indexS{quies} accipiant, duo
\edtext{vero excipientes vice incurrentium eodem}{%
\lemma{vero}%
\Bfootnote{%
\textit{(1)}~incurrentes
\textit{(2)}~excipientes vice incurrentium
\textit{(a)}~(eodem qui hi si nihil impediret mo
\textit{(b)}~eodem~\textit{L}}}
quo hi sine impedimento
\edtext{alias}{\lemma{}\Bfootnote{alias \textit{erg.~L}}}
perrexissent, motu\indexS{motus} moveantur.
\pend
\vspace{0.5\baselineskip}
\pstart
\centering
Demonstratio
\pend
\nobreak
\pstart
\noindent
Sint in recta 
\edtext{centrorum\indexS{recta centrorum gravitatis} (fig.~2) signata sex puncta 1,}{%
\lemma{}%
\Bfootnote{%
centrorum \bars{(fig.~2) \textit{erg.}}
\textit{(1)}~facilitat
\textit{(2)}~signata sex puncta
\bars{aequidistantia \textit{gestr.}} 1,~\textit{L}}}
2, 3, 4, 5, 6 et in ipsis 1 et 2 sint centra \textit{A} et \textit{B}, sed in ipsis 3 et 4 sint centra \textit{L} 
\edtext{et \textit{M}, tendantque \textit{A} motu 1\,2 et B motu aequali et simili 2\,3,}{%
\lemma{et \textit{M},}%
\Bfootnote{%
\textit{(1)}~dico absoluto concursu
\textit{(2)}~tendatque
\textit{(3)}~tendantque \textit{A}
\textit{(a)}~celeritate
\textit{(b)}~ motu 1\,2 et \textit{B} celeritate
\textit{(aa)}~motu
\textit{(bb)}~aequali \bars{et simili \textit{erg.}} 2\,3,~\textit{L}}}
sed quiescant \textit{L} et \textit{M}
\edtext{in 3 et 4,}{\lemma{}\Bfootnote{in 3 et 4 \textit{erg.~L}}}
dico post concursum\indexS{concursus} fore, 
\edtext{ut \textit{A} et \textit{B}}{%
\lemma{ut}%
\Bfootnote{%
\textit{(1)}~\textit{A} qui
\textit{(2)}~\textit{B} et \textit{A}~\textit{L}}}
quiescant in 3 et 4 sed \textit{L} et \textit{M} moveantur motu\indexS{motus post ictum} 4\,5, 5\,6, aequali\indexS{motus aequalis} et conveniente ipsis
\edlabel{35_08_30_040_2a}%
\edtext{}{%
\xxref{35_08_30_040_2a}{35_08_30_040_2b}%
\lemma{1\,2}%
\Bfootnote{%
\textit{(1)}~,
\textit{(2)}~et 2\,3. \lbrack/\rbrack\
\textit{(a)}~Nam \bars{fig.~1 \textit{gestr.}} ponamus primum
\textit{(aa)}~esse
\textit{(bb)}~\textit{A}
\textit{(cc)}~cum \textit{A}, \textit{B}, \textit{C}, \textit{D} sunt in 1, 2, 3, 4
\textit{(aaa)}~esse
\textit{(bbb)}~esse 1\,2 aequ. 3\,4. Cum ergo et
\textit{(b)}~Incipiamus a casu faciliore,
\textit{(c)}~sit distantia ipsorum \textit{A} et \textit{B} ante concursum
\textit{(d)}~Distantia ipsorum \textit{A} et \textit{B} ante concursum sit aequalis
\textit{(e)}~Ponamus \lbrack...\rbrack\ aequalem~\textit{L}}}%
1\,2 et 2\,3.
\pend
\begin{edfigure}%
\includegraphics[scale=1]{%
gesamttex/edit_VIII-4/images/LH_35_08_30_040_d2_040r.pdf}
\edcaption{\lbrack\textit{Fig.~2}\rbrack}
\end{edfigure}
\pstart
Ponamus distantiam ipsorum \textit{A} et \textit{B} ante concursum\indexS{concursus}\indexS{distantia corporum ante concursum} aequalem\edlabel{35_08_30_040_2b}
distantiae ipsorum \textit{L} et \textit{M}. 
\edtext{Et intelligamus momentum quo motus \textit{B}}{%
\lemma{}%
\Bfootnote{Et \bars{in fig.~1 \textit{gestr.}} intelligamus
\textit{(1)}~quand
\textit{(2)}~quo motus~\textit{B}~\textit{L}}}
eo pervenit, ut sit \textit{BL} aequ. \textit{AB} vel 
\edtext{\textit{LM}, hic}{%
\lemma{\textit{LM},}%
\Bfootnote{%
\textit{(1)}~hic esto
\textit{(2)}~hic~\textit{L}}}
\textls{status primus} esto. Sequitur post intervallum temporis 
\edtext{determinatum \textls{status secundus} quo transmisso intervallo punctorum ipse globus \textit{B} ferit \textit{L} in 3\lbrack,\rbrack\ A vero}{%
\lemma{determinatum}%
\Bfootnote{%
\textit{(1)}~tempus
\textit{(2)}~mo
\textit{(3)}~\textls{status secundus} \bars{fig.~2 \textit{gestr.}} quo
\textit{(a)}~\textit{A} quidem attingit punctum 2
\textit{(b)}~transmisso intervallo punctorum
\textit{(aa)}~A attingit punctum
\textit{(bb)}~\bars{ipse globus \textit{erg.}} \textit{B} ferit \textit{L} \bars{in 3 \textit{erg.}}
\textit{(aaa)}~sed \textit{A} tantum abest a
\textit{(bbb)}~sed \textit{A} se collocat in locum ipsius \textit{B}, nempe in 2. Ex hoc \textls{per suppositionem} colligitur status tertius,
\textit{(ccc)}~et per suppositionem
\textit{(ddd)}~sed \textit{A}
\textit{(eee)}~\textit{A} vero~\textit{L}}}
se collocat in locum ipsius \textit{B} nempe in 2.
Sed \textls{per suppositionem} hinc sequitur \textit{L} quidem pelli loco, 
\edtext{sed \textit{B} ejus loco quiescere.}{%
\lemma{sed \textit{B}}%
\Bfootnote{%
\textit{(1)}~in ejus locum
\textit{(2)}~ejus loco quiescere.~\textit{L}}}
Unde \textls{statu} 
\edtext{\textls{tertio} post}{%
\lemma{\textls{tertio}}%
\Bfootnote{%
\textit{(1)}~fit ut \textit{A} 
\textit{(a)}~pergit
\textit{(b)}~pergens
\textit{(2)}~post~\textit{L}}}
aequale priori temporis intervallum \textit{A}
\edtext{quidem pergens feriet \textit{B} quiescentem\indexS{globus quiescens}}{%
\lemma{quidem}%
\Bfootnote{%
\textit{(1)}~ferie
\textit{(2)}~feriet \textit{B} i
\textit{(3)}~pergens feriet \textit{B}
\textit{(aa)}~quiescens
\textit{(bb)}~quiescentem~\textit{L}}}
in 3 sed \textit{L} 
\edtext{in \textit{B} propulsus}{%
\lemma{in \textit{B}}%
\Bfootnote{%
\textit{(1)}~propulsum
\textit{(2)}~propulsus~\textit{L}}}
jam feriet \textit{M} 
\edtext{hactenus quiescentem}{%
\lemma{hactenus}%
\Bfootnote{%
\textit{(1)}~quiescens
\textit{(2)}~quiescentem~\textit{L}}}
in 4. Unde \textls{per suppos.} \textit{A} quidem quietem accipiet
\edtext{in 3 non turbandam,\indexS{quies non turbanda}}{%
\lemma{in 3}%
\Bfootnote{%
\textit{(1)}~non \bars{amplius \textit{gestr.}} turbandam, sed \textit{B},
\textit{(a)}~supr 
\textit{(b)}~motu ejus \bars{accepto \textit{streicht Hrsg.}}
\textit{(aa)}~feriet 
\textit{(bb)}~\textls{statu} jam \textls{quarto} feriet \textit{L} in 4, dum interim \textit{M} ab \textit{L} propulsus pervenit 
\textit{(aaa)}~ad
\textit{(aaaa)}~6
\textit{(bbbb)}~5
\textit{(bbb)}~ad 5. Unde \textls{statu quinto} tandem,
\textit{(aaaa)}~man
\textit{(bbbb)}~etiam \textit{B} \bars{iterum \textit{erg.}} loco pellens \textit{L} (\textls{per suppos.}) 
loco ejus \bars{in 4 \textit{streicht Hrsg.}} quietem accipiet non amplius turbandam;
\textit{(aaaaa)}~et ita \textit{L}
\textit{(bbbbb)}~et \textls{statu quinto} quie
\textit{(2)}~sed \textit{B} post tertium intervallum seu \textls{quarto statu} feriet \textit{L} in 4, dum interea \textit{M} pervenit ad \textlangle5\textrangle\ \bars{et \textit{gestr.}}
\textit{(3)}~non turbandam,~\textit{L}}}
\textit{B} vero recipiet ejus motum; et similiter \textit{L} quiescet in 4, loco \textit{M}, sed \textit{M}
recipiet ejus 
\edtext{motum non turbandum.\indexS{motus non turbandus} Quare jam post}{%
\lemma{motum}%
\Bfootnote{%
\textit{(1)}~.~Unde
\textit{(2)}~\bars{non turbandum \textit{erg.}}. Quare jam
\textit{(a)}~\textls{statu quarto}
\textit{(aa)}~\textit{A} manente in 3, \textit{B} feriet \textit{L} in 4, et \textit{M}
\textit{(bb)}~post
\textit{(b)}~post~\textit{L}}}
tertium intervallum \textls{statu quarto} \textit{A} manente in 3, \textit{B} feriet \textit{L} in 4,
et \textit{M} perveniet 
\edtext{ad 5. Et}{%
\lemma{ad 5.}%
\Bfootnote{%
\textit{(1)}~Ita
\textit{(2)}~Et~\textit{L}}}
\textls{per suppos.} \textit{B} quietem incipiet in 4 non amplius 
\edtext{turbandam, sed}{%
\lemma{turbandam,}%
\Bfootnote{%
\textit{(1)}~et
\textit{(2)}~sed~\textit{L}}}
\textit{L} accipiet ejus motum
\edtext{non amplius turbandum.\indexS{motus non turbandus}}{\lemma{}\Bfootnote{non amplius turbandum \textit{erg.~L}}}
Ergo denique post quartum intervallum \textls{statu quinto} \textit{A} et \textit{B} quiescentibus
in 3 et 4, \textit{L} perveniet in 5 motu 4\,5, et \textit{M} in 6 motu 5\,6. Ut asseverabatur.
\edtext{Quod si intervallum
\href{\Scan{LH_35_08_30_040v}}%
{\lbrack40~v\textsuperscript{o}\rbrack}\ 
ipsorum \textit{A} et \textit{B} ante concursum}{%
\lemma{Quod si}%
\Bfootnote{%
\textit{(1)}~intervalla
\href{\Scan{LH_35_08_30_040v}}%
{\lbrack40~v\textsuperscript{o}\rbrack}\ 
ipsorum \textit{A} et \textit{B}, non sint eadem
\textit{(2)}~intervallum 
\href{\Scan{LH_35_08_30_040v}}%
{\lbrack40~v\textsuperscript{o}\rbrack}\ 
\lbrack...\rbrack\ concursum~\textit{L}}}
non sit idem quod ipsorum \textit{L} et \textit{M}, patet nihil mutari in motibus,
sed tantum in intervallis temporum, semper enim necesse est
\textit{A} et \textit{B} redigi ad quietem sed \textit{L} et \textit{M} postremo moveri, 
nullum alium vero accipere, quam qui erat ipsorum \textit{A} et \textit{B}.
\pend
\pstart
\centering
Theorema 2.
\pend
\pstart
\noindent
Idem est iisdem positis, si plures sint 
\edtext{globi,\indexS{globus quiescens} modo quiescentium et in}{%
\lemma{globi}%
\Bfootnote{%
\textit{(1)}~incipientes
\textit{(2)}~,~modo
\textit{(a)}~ex
\textit{(b)}~quiescentium et
\textit{(aa)}~mot
\textit{(bb)}~in~\textit{L}}}
motu existentium\indexS{globus in motu existens} numerus sit aequalis.
\pend
\pstart
Sed operae pretium foret idem distincte ponere in tribus.
\pend
\pstart
\edtext{Fortasse magis commodum erit demonstrare, si
\edtext{unus globus\indexS{globus in quotcunque globos quiescentes incurrens}}{%
\lemma{unus}%
\Bfootnote{%
\textit{(1)}~modo
\textit{(2)}~globus~\textit{L}}}
incurrat in quotcunque aequales et quiescentes\indexS{globus quiescens}
\edtext{unum solum}{%
\lemma{unum}%
\Bfootnote{%
\textit{(1)}~quiescere
\textit{(2)}~ultimum so
\textit{(3)}~solum~\textit{L}}}
maxime remotum progredi, caeteros quiescere una cum advena.}{%
\lemma{Fortasse \bracks{...} advena}%
\applabel{FnC:LH_35_08_30_040-advena}
\Cfootnote{Die hier skizzierte Untersuchung leistet Leibniz tatsächlich im Konzept N.\;\ref{58297} (August 1697 bis April 1703).}}
Cujus rei manifesta ratio est, quid 
\edtext{quilibet incurrens in}{%
\lemma{quilibet}%
\Bfootnote{%
\textit{(1)}~praecedentem quiescen
\textit{(2)}~incurrens in~\textit{L}}}
anteriorem quiescentem\indexS{globus quiescens} redigitur ad quietem, et suum illi tribuit motum.\indexS{globus alteri motum tribuens}
Jam si unus sit incurrens, statim
\edtext{redigitur ad quietem,}{%
\lemma{redigitur}%
\Bfootnote{%
\textit{(1)}~et qui
\textit{(2)}~ad quietem,~\textit{L}}}
et de caetero res redit, 
\edtext{ac si esset unus}{%
\lemma{ac si}%
\Bfootnote{%
\textit{(1)}~alter
\textit{(2)}~esset
\textit{(a)}~incurren
\textit{(b)}~ex
\textit{(c)}~unus~\textit{L}}}
incurrens excipientium numerus uno minor.
Ergo tandem res eo redit ac si esset excipiens\indexS{globus excipiens} unus, et incurrens\indexS{globus incurrens} unus, id est, unus movebitur, qui est excipiens ultimus.
\pend
\pstart
Si duo\indexS{globi duo directe incurrentes in plures quiescentes} incurrant
\edtext{in alios, primo incurrente, fient}{%
\lemma{in}%
\Bfootnote{%
\textit{(1)}~duos
\textit{(2)}~alios,
\textit{(a)}~ultimo incurre
\textit{(b)}~primo incurrente
\textit{(aa)}~, res
\textit{(bb)}~solo, 
\textit{(aaa)}~res eo
\textit{(bbb)}~abstrahamus a secundo, et
\textit{(aaaa)}~ut
\textit{(bbbb)}~fiet ut
\textit{(ccc)}~fient~\textit{L}}}
omnia qua propositione praecedente, 
\edtext{nempe omnes quiescent, et solus}{%
\lemma{nempe}%
\Bfootnote{%
\textit{(1)}~ipse
\textit{(2)}~omnes quiescent, et
\textit{(a)}~ultimus solus
\textit{(b)}~solus~\textit{L}}}
remotissimus movebitur unde secundo superveniente, idem est, quod ante, 
seu rediimus ad statum priorem.
Ergo rursus movebitur 
\edtext{remotissimus\bracks{,} quiescent}{%
\lemma{remotissimus}%
\Bfootnote{%
\textit{(1)}~secundi
\textit{(2)}~ea
\textit{(3)}~quiescent~\textit{L}}}
caeteri. Et idem est 
\edtext{si sint}{%
\lemma{}%
\Bfootnote{%
si \bars{si \textit{streicht Hrsg.}} sint~\textit{L}}}
plures.
\pend
\pstart
Nec refert quod 
\edtext{secundus ferire}{%
\lemma{secundus}%
\Bfootnote{%
\textit{(1)}~feriat ante
\textit{(2)}~ferire~\textit{L}}}
fortasse possit ante quam prioris operatio sit transacta.
\edtext{Nam saltem secundi}{%
\lemma{Nam}%
\Bfootnote{%
\textit{(1)}~saltem
\textit{(2)}~saltem
\textit{(a)}~prioris operatio non p
\textit{(b)}~secundi~\textit{L}}}
operatio non perveniet ad remotiores, quam omnia primi ibi jam sint transacta. Unde res eodem redit.
\pend
\pstart
Nihil etiam refert, etsi plures sint incurrentes\indexS{globi plures directe incurrentes} quam excipientes,\indexS{globus excipiens}
\edtext{nam}{\lemma{}\Bfootnote{nam \textit{erg.~L}}}
semper res eo 
\edtext{redit ope}{%
\lemma{redit}%
\Bfootnote{%
\textit{(1)}~,~ac si incurrentes null
\textit{(2)}~ope~\textit{L}}}
vicariorum, ac si incurrentes\indexS{globus incurrens} obstaculum nullum sensissent.
\pend
\pstart
Superest jam ut ostendamus idem fieri, et si utrinque sit motus, si 
\edtext{modo duobus concurrentibus\indexS{globi duo concurrentes} ponamus fieri permutationem.\indexS{permutatio motus in concursu}}{%
\lemma{modo}%
\Bfootnote{%
\textit{(1)}~in consul
\textit{(2)}~duobus concurrentibus ponamus
\textit{(a)}~esse
\textit{(b)}~fieri
\textit{(aa)}~decussa
\textit{(bb)}~permutationem.~\textit{L}}}
\pend
\pstart
Nam cum 
\edtext{semper in concursu\indexS{concursus aequalium} aequalium}{%
\lemma{semper}%
\Bfootnote{%
\textit{(1)}~omnia fiant ut ante per motum
\textit{(2)}~in concursu aequalium~\textit{L}}}
medio motu omnia fiant perinde ac si esset penetratio\indexS{penetratio quasi dimensionum in concursu corporum} dimensionorum, postremo etiam idem 
\edlabel{35_08_30_040_3a}%
\edtext{}{%
\xxref{35_08_30_040_3a}{35_08_30_040_3b}%
\lemma{fiet.}%
\Bfootnote{%
\textit{(1)}~Ni
\textit{(2)}~\bars{Cum ergo non sit nisi concursu aequalium \textit{streicht Hrsg.}} fiet
\textit{(3)}~Hinc~\textit{L}}}%
fiet.
\pend
\pstart
Hinc\edlabel{35_08_30_040_3b}
infero aliquid pulcherrimum, etsi plures globi simul diversis motibus moveantur in eadem recta,\indexS{globi in recta diversis motis moti} 
post concursum\indexS{concursus} perinde faciendum, ac si per penetrationem\indexS{penetratio dimensionum} dimensionum
quilibet inoffensus progressus\indexS{progressus corporis inoffensus} fuisset\bracks{,} seu perinde ac si 
\edtext{fuissent singuli}{%
\lemma{fuissent}%
\Bfootnote{%
\textit{(1)}~in eadem linea omne
\textit{(2)}~singuli~\textit{L}}}
in linea sua priori
\edtext{parallela seu}{%
\lemma{parallela}%
\Bfootnote{%
\textit{(1)}~umbraeque solares in
\textit{(2)}~seu~\textit{L}}}
prorsus ut eorum 
\edtext{umbrae solares\indexS{umbra solaris in pariete}}{%
\lemma{umbrae}%
\Bfootnote{%
\textit{(1)}~in
\textit{(2)}~solares~\textit{L}}}
in pariete moverentur nam umbrae se impedire non possunt.
\pend
\pstart
Idem est si oblique concurrant,\indexS{concursus obliquus centralis} sed tamen centraliter 
\edtext{obliquitatem ascribendo}{%
\lemma{obliquitatem}%
\Bfootnote{%
\textit{(1)}~enim
\textit{(2)}~ascribendo~\textit{L}}}
motui cuidam transverso irrefracto.
\pend
\begin{edfigure}%
\includegraphics[scale=1]{%
gesamttex/edit_VIII-4/images/LH_35_08_30_040_d3_040v.pdf}
\edcaption{%
\edtext{\lbrack\textit{Fig.~3}\rbrack}{%
\lemma{\lbrack\textit{Fig.~3}\rbrack}%
\Cfootnote{%
Zwei gestrichene Entwürfe zum Diagramm werden nicht wiedergegeben.}}%
}
\end{edfigure}

\pstart
Sed si globi sint inaequales, vel si aequales connexi\indexS{globi aequales connexi} ponantur, ut a seinvicem divelli
\edtext{nequeant, sitque}{%
\lemma{nequeant,}%
\Bfootnote{%
\textit{(1)}~hic quam
\textit{(2)}~fit
\textit{(3)}~sitque~\textit{L}}}
vis Elastica\indexS{vis Elastica} eos 
connectens, considerandum est \textit{B} incurrens in \textit{L} dare ipsi motum suum.
Sed \textit{L} tendens versus 
\edtext{\textit{M} cogi intendere Elastrum\indexS{Elastrum} interceptum,}{%
\lemma{\textit{M}}%
\Bfootnote{%
\textit{(1)}~cogere
\textit{(2)}~cogi intendere Elastrum \bars{suum \textit{gestr.}} interceptum,~\textit{L}}}
quo facto Elastrum\indexS{Elastrum} aget in \textit{M}, eique dabit aliquid motus sed parvi,
sed quia \textit{M} tardius longe cedit, quam \textit{L} insequitur, continue
magis intendetur Elastrum,\indexS{Elastrum continue magis intensum} donec vis\indexS{vis} ipsius \textit{L} tota
partim Elastro partim ipsi \textit{M} sit tributa,\indexS{translatio virium} quo facto vi Elastri\indexS{vis Elastri} rursus a se recedent, et \textit{L} quidem
redibit ad locum priorem ultraque tendens residuo impetu feriet 
\edtext{\textit{B}, qui}{%
\lemma{\textit{B},}%
\Bfootnote{%
\textit{(1)}~quod
\textit{(2)}~qui~\textit{L}}}
globus proinde reflectetur sed \textit{L} deinde trahetur ab ipso 
\edtext{\textit{M} et}{%
\lemma{\textit{M}}%
\Bfootnote{%
\textit{(1)}~donec
\textit{(2)}~Elastro
\textit{(3)}~et~\textit{L}}}
elastro in contrarium intenso,\indexS{Elastrum intensum} seu distenso.\indexS{Elastrum distensum}
\edtext{Quae omnia distinctius exponi merentur.}{%
\applabel{FnC:LH_35_08_30_040-Papin}%
\lemma{Quae omnia distinctius exponi merentur}%
\Cfootnote{%
Leibniz führte im Briefwechsel mit \indexN{\textls{Papin} (Papinus), Denis 1647\textendash1713}Papin zwischen Juli 1696 und Januar 1698 eine rege Debatte über das Gedankenexperiment einer zwischen den stoßenden Körpern befindlichen Feder: Siehe dazu die Datierung zu N.\;\ref{60360} und die dort zitierten Briefe.}}
\pend
\begin{edfigure}%
\includegraphics[scale=1]{%
gesamttex/edit_VIII-4/images/LH_35_08_30_040_d4_040v.pdf}
\edcaption{%
\edtext{\lbrack\textit{Fig.~4}\rbrack}{%
\lemma{\lbrack\textit{Fig.~4}\rbrack}%
\Cfootnote{Das Diagramm ist wie auch der nachfolgende Text eine Ergänzung.}}%
}
\end{edfigure}
\pstart
\edtext{Quid si globus\indexS{globus in duos simul incurrens} \textit{A} incurrat in \textit{B} et \textit{C} simul, qui tangunt eodem modo \textit{D} quo \textit{A}\bracks{?}
An \textit{A} quiescet et ejus loco \textit{D} perget? Non puto.}{%
\lemma{}%
\Bfootnote{%
Quid \lbrack...\rbrack\ puto. \textit{erg.~L}}}
\pend
